\subsection{Training Data}
\label{sec:training_data}

Training data were drawn from the RGB color cube, which encompasses all colors representable in the RGB color model. Each sample is represented as a 3D vector in RGB space, with each channel in $[0,1]$. The cube was subdivided uniformly along each axis, and training samples were generated at the grid points. In our experiments, we used a $8 \times 8 \times 8$ grid, yielding $512$ training samples, shown in~\cref{fig:rgb-cube}.

\begin{figure}[htb]
    \centering
    \begin{subfigure}[b]{0.32\textwidth}
        \centering
        \includegraphics[width=\textwidth]{figures/ex-1.1.1-rgb-cube-front.png}
        \caption{Front}
    \end{subfigure}
    % \hfill
    \hspace{1em}
    \begin{subfigure}[b]{0.32\textwidth}
        \centering
        \includegraphics[width=\textwidth]{figures/ex-1.1.1-rgb-cube-back.png}
        \caption{Back}
    \end{subfigure}
    \caption{\textbf{The RGB cube as training data.} Two views of the cube are shown, both oriented such that the black-to-white diagonal runs from bottom-to-top; thus red, blue, and green are nearer the bottom, whereas cyan, yellow, and magenta are nearer the top. Grays are located in the center of the cube (not visible).
    \textit{a}:~View facing the warm hues, with red in the middle and yellow and magenta on either side. \textit{b}:~View facing the cool hues, with cyan in the middle and blue and green on either side.
    }
    \label{fig:rgb-cube}
\end{figure}

Superficially, the cube resembles some of our latent space plots---however the spaces are quite different: the RGB cube is solid (the magnitudes of $\vx$ and $\vy$ are significant), whereas our latent spaces are hyperspherical ($\hat{\vz}$ is purely directional).
